% !TeX root = ../main.tex
% Fonte: Lezione 7.pdf, pp. 10--18.
\section{Funzioni definite a tratti}

\subsection{Parte positiva e parte negativa}

Per ogni reale $x$ si definiscono la \emph{parte positiva} e la \emph{parte negativa} mediante
\[
x^+=\max\{x,0\}=\begin{cases}x,&x\ge0,\\0,&x<0,\end{cases}
\qquad
x^-=-\min\{x,0\}=\begin{cases}-x,&x\le0,\\0,&x>0.\end{cases}
\]
\textbf{Segno e identità.} Entrambe sono non negative: la parte negativa misura il modulo della componente negativa e non è un numero negativo. Distinguendo i casi $x\ge0$ e $x<0$ si verificano le identità
\[
\boxed{x=x^+-x^-,\qquad |x|=x^++x^-}.
\]


\textbf{Parti di una funzione.} Le stesse definizioni si applicano punto per punto a una funzione: $f^+(x)=\max\{f(x),0\}$ e $f^-(x)=\max\{-f(x),0\}$, con $f=f^+-f^-$ e $|f|=f^++f^-$.

\begin{schema}
\begin{tikzpicture}[x=1cm,y=.85cm]
  \draw[asse] (-2.4,0) -- (2.5,0) node[right] {$x$};
  \draw[asse] (0,-.3) -- (0,2.6) node[above] {$y$};
  \draw[intervallo] (-2.1,0) -- (0,0) -- (2.1,2.1);
  \node[above left] at (2.1,2.1) {$x^+$};
\end{tikzpicture}
\hspace{1cm}
\begin{tikzpicture}[x=1cm,y=.85cm]
  \draw[asse] (-2.4,0) -- (2.5,0) node[right] {$x$};
  \draw[asse] (0,-.3) -- (0,2.6) node[above] {$y$};
  \draw[intervallo] (-2.1,2.1) -- (0,0) -- (2.1,0);
  \node[above right] at (-2.1,2.1) {$x^-$};
\end{tikzpicture}
\end{schema}

\subsection{Parte intera e parte frazionaria}
\enlargethispage{2pt}

\begin{definizione}
La \emph{parte intera} di $x\in\R$, o \emph{floor}, è l'unico intero $m$ tale che $m\le x<m+1$. Si indica con $\lfloor x\rfloor$.\footnote{Si usano anche $[x]$ per la parte intera e $(x)$ per la parte frazionaria; qui preferiamo $\lfloor x\rfloor$ e $\{x\}$ per distinguere chiaramente le due funzioni dalle altre parentesi.} La \emph{parte frazionaria} è $\{x\}=x-\lfloor x\rfloor$.
\end{definizione}

\begin{proof}[Esistenza e unicità]
\textbf{Esistenza.} Archimede e il principio del minimo intero forniscono, come nella prova della densità dei razionali, un intero $m$ con $x-1<m\le x$: basta applicare a $x-1$ il fatto che per ogni reale $t$ esiste un intero in $(t,t+1]$. Si ottiene $m\le x<m+1$.

\textbf{Unicità.} Se due interi $m<k$ soddisfacessero entrambe le condizioni, si avrebbe $m+1\le k\le x<m+1$, assurdo: due interi distinti distano almeno $1$.
\end{proof}

Per ogni reale $x$ si ha $0\le\{x\}<1$; dunque
\[
\lfloor\,\cdot\,\rfloor\colon\R\to\Z,\qquad
\{\,\cdot\,\}\colon\R\to[0,1).
\]
\textbf{Grafici.} Per $m\le x<m+1$, il grafico della parte intera è il segmento orizzontale $y=m$ e quello della parte frazionaria è il segmento $y=x-m$, di pendenza $1$. In entrambi, l'estremo sinistro è incluso e il destro escluso.

\begin{proof}[Periodo principale]
\textbf{Periodo $1$ della parte frazionaria.} Se $m=\lfloor x\rfloor$, allora $m+1\le x+1<m+2$: dunque $\lfloor x+1\rfloor=\lfloor x\rfloor+1$ e $\{x+1\}=\{x\}$.

\textbf{Minimalità del periodo.} Ogni periodo $T>0$ soddisfa $\{T\}=\{0\}=0$, quindi $T\in\N^*$ e $T\ge1$. Il periodo $1$ è pertanto il periodo principale.
\end{proof}

\begin{esempio}
Si ha $\lfloor2{,}7\rfloor=2$ e $\{2{,}7\}=0{,}7$, ma $\lfloor-1{,}3\rfloor=-2$ e $\{-1{,}3\}=0{,}7$. La parte intera arrotonda verso il basso, anche per i negativi; non coincide con l'eliminazione delle cifre dopo la virgola.
\end{esempio}

\begin{schema}
\begin{tikzpicture}[x=.8cm,y=.65cm]
  \draw[asse] (-2.3,0) -- (3.4,0) node[right] {$x$};
  \draw[asse] (0,-2.4) -- (0,3.4) node[above] {$y$};
  \foreach \m in {-2,-1,0,1,2} {
    \draw[intervallo] (\m,\m) -- ({\m+1},\m);
    \fill[coloreschema] (\m,\m) circle (1.7pt);
    \draw[estremo] ({\m+1},\m) circle (1.7pt);
  }
  \draw[intervallo] (3,3) -- (3.25,3);
  \fill[coloreschema] (3,3) circle (1.7pt);
  \foreach \m in {-2,-1,1,2,3} \node[below=4pt] at (\m,0) {$\m$};
  \node[below] at (.5,-2.9) {$y=\lfloor x\rfloor$};
\end{tikzpicture}
\hspace{.6cm}
\begin{tikzpicture}[x=.8cm,y=1.3cm]
  \draw[asse] (-2.3,0) -- (3.4,0) node[right] {$x$};
  \draw[asse] (0,-.4) -- (0,1.6) node[above] {$y$};
  \draw[dotted,black!40] (-2.1,1) -- (3.1,1);
  \foreach \m in {-2,-1,0,1,2} {
    \draw[intervallo] (\m,0) -- ({\m+1},1);
    \fill[coloreschema] (\m,0) circle (1.7pt);
    \draw[estremo] ({\m+1},1) circle (1.7pt);
  }
  \draw[intervallo] (3,0) -- (3.25,.25);
  \fill[coloreschema] (3,0) circle (1.7pt);
  \foreach \m in {-2,-1,1,2,3} \node[below=4pt] at (\m,0) {$\m$};
  \node[above left] at (0,1) {$1$};
  \node[below] at (.5,-.8) {$y=\{x\}$};
\end{tikzpicture}
\end{schema}

\section{Trasformazioni dei grafici}
% Fonte: Lezione 7.pdf, p. 13.

Sia $D$ il dominio di $f$. Dal grafico di $y=f(x)$ si possono ottenere altri grafici seguendo come cambia un suo punto $(x,f(x))$:
\begin{itemize}
  \item \textbf{Traslazione verticale.} $y=f(x)+a$: ogni punto diventa $(x,f(x)+a)$; il grafico si trasla verticalmente di $a$, verso l'alto se $a>0$ e verso il basso se $a<0$;
  \item \textbf{Traslazione orizzontale.} $y=f(x-b)$: ogni punto $(u,f(u))$ diventa $(u+b,f(u))$; la traslazione è orizzontale, verso destra se $b>0$ e verso sinistra se $b<0$. Il dominio diventa $D+b=\{u+b:u\in D\}$;
  \item \textbf{Riflessione rispetto all’asse $x$.} $y=-f(x)$: ogni punto diventa $(x,-f(x))$; il grafico viene riflesso rispetto all'asse $x$.
\end{itemize}
Per $b>0$, il segno meno \emph{dentro} $f(x-b)$ corrisponde a uno spostamento verso destra: per conservare l'argomento $u$ si prende $x=u+b$. Il verso della traslazione dipende quindi dal segno di $b$.

\begin{schema}
\begin{tikzpicture}[x=.8cm,y=.55cm]
  \draw[asse] (-.7,0) -- (3.5,0) node[right] {$x$};
  \draw[asse] (0,-2.3) -- (0,4.2) node[above] {$y$};
  \draw[black!45,domain=-.4:3.05,samples=90] plot (\x,{\x*(\x-2)^2});
  \draw[intervallo,domain=-.4:3.05,samples=90] plot (\x,{\x*(\x-2)^2+1});
  \node[below] at (1.3,-2.9) {$f(x)+1$};
\end{tikzpicture}
\hspace{.35cm}
\begin{tikzpicture}[x=.8cm,y=.55cm]
  \draw[asse] (-.7,0) -- (4.5,0) node[right] {$x$};
  \draw[asse] (0,-2.3) -- (0,4.2) node[above] {$y$};
  \draw[black!45,domain=-.4:3.05,samples=90] plot (\x,{\x*(\x-2)^2});
  \draw[intervallo,domain=.6:4.05,samples=90] plot (\x,{(\x-1)*(\x-3)^2});
  \node[below] at (1.8,-2.9) {$f(x-1)$};
\end{tikzpicture}
\hspace{.35cm}
\begin{tikzpicture}[x=.8cm,y=.55cm]
  \draw[asse] (-.7,0) -- (3.5,0) node[right] {$x$};
  \draw[asse] (0,-3.6) -- (0,3.5) node[above] {$y$};
  \draw[black!45,domain=-.4:3.05,samples=90] plot (\x,{\x*(\x-2)^2});
  \draw[intervallo,domain=-.4:3.05,samples=90] plot (\x,{-\x*(\x-2)^2});
  \node[below] at (1.3,-4.1) {$-f(x)$};
\end{tikzpicture}

In grigio il grafico iniziale; in colore la trasformazione indicata.
\end{schema}

\Needspace{16\baselineskip}
\section{Un esercizio su dominio, monotonia e inversa}
% Fonte: Lezione 7.pdf, pp. 14--18. Corretto il titolo «non decr.»
% della p. 16, incompatibile con il calcolo decrescente svolto sotto.

\begin{esercizio}
Studiamo la funzione
\[
f(x)=4\arcsin\!\bigl(1-\ln(x-1)\bigr).
\]
Cerchiamo il dominio naturale, il verso di monotonia e l'insieme dei valori $y\in\R$ per i quali $f(x)=y$ ammette una soluzione.
\end{esercizio}

\subsection{Determinazione del dominio}

\begin{svolgimento}
\textbf{Condizioni di definizione.} Il logaritmo richiede $x-1>0$; l'arcoseno richiede che il proprio argomento sia in $[-1,1]$. Quindi
\begin{align*}
x\in D
&\ \Longleftrightarrow\
\begin{cases}x>1,\\-1\le1-\ln(x-1)\le1\end{cases}\\
&\ \Longleftrightarrow\
\begin{cases}x>1,\\0\le\ln(x-1)\le2\end{cases}\\
&\ \Longleftrightarrow\
\begin{cases}x>1,\\1\le x-1\le e^2.\end{cases}
\end{align*}
L'ultima doppia disuguaglianza equivale a $2\le x\le1+e^2$ e implica già $x>1$. Pertanto
\[
\boxed{D=[2,1+e^2]}.
\]
\end{svolgimento}

\subsection{Monotonia}

\begin{svolgimento}
Siano $x_1,x_2\in D$ con $x_1<x_2$. Il logaritmo è strettamente crescente, perciò
\[
x_1-1<x_2-1
\ \Longrightarrow\ \ln(x_1-1)<\ln(x_2-1)
\ \Longrightarrow\ 1-\ln(x_1-1)>1-\ln(x_2-1).
\]
L'arcoseno è strettamente crescente e moltiplicare per $4>0$ conserva il verso: ne segue $f(x_1)>f(x_2)$. Dunque $f$ è \emph{strettamente decrescente}, quindi iniettiva sul proprio dominio.

\end{svolgimento}

\Needspace{30\baselineskip}

\subsection{Immagine, equazione e inversa}

\begin{svolgimento}
\textbf{Necessità.} Poiché l'arcoseno assume valori in $[-\pi/2,\pi/2]$, una condizione necessaria per $f(x)=y$ è $-2\pi\le y\le2\pi$.

\textbf{Sufficienza: costruzione della soluzione.} Sotto questa condizione $y/4$ appartiene proprio all'intervallo in cui seno e arcoseno sono inversi, e possiamo procedere mediante equivalenze:
\begin{align*}
4\arcsin\!\bigl(1-\ln(x-1)\bigr)=y
&\ \Longleftrightarrow\ 1-\ln(x-1)=\sin(y/4)\\
&\ \Longleftrightarrow\ \ln(x-1)=1-\sin(y/4)\\
&\ \Longleftrightarrow\ x=1+e^{\,1-\sin(y/4)}.
\end{align*}
Per ogni $y\in[-2\pi,2\pi]$ il valore $1-\sin(y/4)$ è in $[0,2]$; quindi la soluzione trovata appartiene a $[2,1+e^2]$ ed è valida. Per i valori di $y$ esterni all'intervallo non ci sono soluzioni. In conclusione,
\[
\operatorname{Im}(f)=[-2\pi,2\pi],\qquad
f^{-1}(y)=1+e^{\,1-\sin(y/4)}\quad(-2\pi\le y\le2\pi).
\]
\textbf{Unicità.} La soluzione è unica per l'iniettività.
\end{svolgimento}

\textbf{Massimo e minimo.} Gli estremi sono raggiunti: $f(2)=2\pi$ è il massimo e $f(1+e^2)=-2\pi$ è il minimo.

\begin{schema}
\begin{tikzpicture}[x=.85cm,y=.4cm]
  \draw[asse] (-.4,0) -- (9.2,0) node[right] {$x$};
  \draw[asse] (0,-7.5) -- (0,7.6) node[above] {$y$};
  \draw[intervallo,domain=-pi/2:pi/2,samples=100]
    plot ({1+exp(1-sin(\x r))},{4*\x});
  \draw[dotted] (2,0) -- (2,{2*pi}) -- (0,{2*pi});
  \draw[dotted] ({1+exp(2)},0) -- ({1+exp(2)},{-2*pi}) -- (0,{-2*pi});
  \fill[coloreschema] (2,{2*pi}) circle (1.8pt);
  \fill[coloreschema] ({1+exp(2)},{-2*pi}) circle (1.8pt);
  \node[below] at (2,0) {$2$};
  \node[above] at ({1+exp(2)},0) {$1+e^2$};
  \node[left] at (0,{2*pi}) {$2\pi$};
  \node[left] at (0,{-2*pi}) {$-2\pi$};
\end{tikzpicture}

Il grafico collega gli estremi inclusi $(2,2\pi)$ e $(1+e^2,-2\pi)$.
\end{schema}
